\section{Введение}
\label{sec:introduction}

\subsection{Эмпатия как поле: сто лет разрозненных прикосновений}

Начиная с ранних психоаналитиков \citep{ferenczi1933, klein1946} и до современной когнитивной нейронауки \citep{decety2011, singer2014}, психология эмпатии накопила несколько крупных, внутренне богатых, но плохо стыкующихся между собой линий исследования. Психоаналитическая традиция описывает расщепление, вытеснение и защитные механизмы, превращающие переживание другого в свою противоположность \citep{freud1915, fairbairn1952, bion1962, fonagy2002}. Теория привязанности разрабатывает средовые предшественники --- качество ранних отношений как предиктор способности к аффективному резонансу в зрелом возрасте \citep{bowlby1969, main1986, lyonsruth1999}. Когнитивная нейронаука разделяет когнитивную эмпатию (theory of mind) от аффективной (resonance), указывая на различные нейронные субстраты \citep{shamaytsoory2011, decety2011}. Зеркально-нейронная литература предлагает биологический механизм аффективного резонанса \citep{rizzolatti2004, iacoboni2009}. Клиническая традиция психопатии, от \citet{cleckley1941} до \citet{hare1991} и \citet{blair2013}, работает с аффективным дефицитом как структурным свойством некоторой подгруппы индивидов. Созерцательные традиции, от Будды до современной non-dual литературы \citep{godman1985, dunne2011, josipovic2014}, описывают иной слой сознания --- наблюдающее отношение, мета-осознанность, witnessing, --- который, согласно этим традициям, модулирует все остальные психические функции.

Каждая из этих линий развивалась независимо, со своими инструментами, со своей лексикой, со своими validation-стандартами. Это не теоретический упрёк: детальная работа внутри каждой парадигмы приносит знания, которые были бы невозможны в эклектическом объединении. Но это эмпирическая проблема: исследователь, который хочет измерить эмпатию у конкретного человека и сделать обоснованные заявления о её ресурсах и уязвимостях, оказывается в ситуации, где у него есть пять или шесть частичных карт одной территории, и ни одна из них не покрывает всю территорию.

В данной статье мы предлагаем не седьмую карту поверх шести существующих. Мы предлагаем \emph{измерительный каркас}, в котором конструкты, измеримые внутри каждой из существующих традиций, получают общую количественную координатную систему, позволяющую сравнивать и соотносить их в одной выборке. Мы называем этот каркас \DEIC{} (\textit{Dynamic Empathy and Inner Conflict} --- условное название, подчёркивающее фокус на взаимодействии эмпатии и защитных механизмов). Психометрически \DEIC{} --- шестифакторная модель с одним дополнительным формативным конструктом (вниманием, A), проверенная конфирматорным анализом на выборке $N = 1426$.

\subsection{Эмпатия как дефолт: перевёрнутая рамка измерения}
\label{subsec:empathy_as_default}

Прежде чем описывать конкретные факторы модели, сформулируем инвертирующий тезис, на котором \DEIC{} построен и который отличает его от большинства существующих опросников эмпатии. Стандартная традиция измерения --- от IRI \citep[Interpersonal Reactivity Index;][]{davis1983} через Empathy Quotient \citep{baroncohen2011} до современных многомерных инструментов --- исходит из неявной предпосылки, что эмпатия есть индивидуальная способность, переменная по величине: у одних людей её больше, у других меньше, задача инструмента --- количественно эту величину оценить. Мы исходим из противоположной предпосылки.

\textbf{Эмпатия --- дефолтное состояние человеческой психики.} Аффективный резонанс с чужим состоянием --- не достижение, не ресурс, не черта, которую одни развили, а другие нет; это базовая функция социального мозга приматов \citep{rizzolatti2004, iacoboni2009, decety2011}, неврологически поддержанная системой зеркальных нейронов и аффективного резонанса, которая у младенцев работает с первых месяцев жизни. \emph{Наблюдаемые различия между людьми по эмпатии --- это не различия по её количеству, а различия по механизмам её подавления.} Модель предлагает поэтому измерять не эмпатию как таковую, а \emph{формы и силу не-эмпатии} --- активные и пассивные процессы, блокирующие базовую способность.

Из этого переворота следует разделение, которое в \DEIC{} является центральным концептуальным инструментом. Существуют два архитектурно различных канала не-эмпатии.

\textbf{Пассивная блокировка.} Эмпатический отклик начинается, но блокируется защитной системой, потому что само переживание --- \emph{неважно, чьё, чужое или собственное}, --- невыносимо для психики в её текущей конфигурации. Здесь необходимо сразу снять одно распространённое упрощение: пассивная блокировка не есть блокировка «чувствования другого» отдельно от «чувствования себя». Это один и тот же канал. Способность вынести аффективный резонанс с чужой болью и способность вынести собственный внутренний материал --- структурно одна функция, а не две разные: реверсивная чувствительность контейнирующей системы, которая либо держит аффект, либо не держит. Психоаналитическая и mentalization-традиция фиксирует это явно: $\alpha$-функция у \citet{bion1962} --- способность метаболизировать эмоциональный опыт --- работает симметрично на собственных и воспринятых состояниях; reflective function у \citet{fonagy2002} по определению двунаправленна (способность ментализировать других предиктируется способностью ментализировать себя и наоборот); holding у \citet{winnicott1965} возможно у матери ровно в той мере, в какой она выдерживает собственные состояния. Нейробиологическая линия \citet{schore2001} показывает, что субстрат self-regulation и other-attunement --- общий: правополушарные сети, отвечающие за регуляцию собственного возбуждения, суть те же сети, которые обеспечивают аффективную настройку на другого. Поэтому человек в режиме пассивной блокировки \emph{не может} чувствовать в обе стороны одновременно: он одинаково отрезан от слёз над чужой потерей и от собственной усталости, одинаково не замечает растущей тревоги в партнёре и в себе самом. Феноменологически это проявляется как онемение, диссоциация, аффективное уплощение, внутренний конфликт «должен чувствовать, но не получается» --- и важно, что эта «должность» относится и к отклику на другого, и к контакту со своим содержанием. Психометрическая координата --- фактор ID (Internal Dissonance). Исторические описания: \citet{ferenczi1933} про «идентификацию с агрессором», \citet{klein1946} про расщепление, \citet{bowlby1969} и \citet{main1986} про дезорганизованные модели привязанности, клиническая диссоциативная литература \citep{putnam1997, vanderkolk2014}.

\textbf{Активное отключение.} Эмпатический канал сохранен и функционально доступен, но не запускается автоматически; он включается или не включается в зависимости от контекстной пригодности и морального нарратива о ситуации. Человек в этом режиме \emph{может} чувствовать выборочно и может \emph{не} чувствовать, когда это удобно. Феноменологически это не онемение, а контроль --- плавная переключаемость. Психометрическая координата --- фактор P (психопатический паттерн). Исторические описания: \citet{cleckley1941}, \citet{karpman1941}, клиническая традиция \citep{hare1991, lilienfeld1990}, нейронаучная линия Блэра \citep{blair1995, blair2013}.

Эти два канала независимы. Они \emph{не} суть «одно и то же в разной пропорции» и не образуют единого спектра. Человек может иметь пассивную блокировку при полном отсутствии психопатического нарратива (типичный клинический профиль депрессии и комплексной травмы); может иметь психопатическую условную архитектуру при отсутствии ID (первичный паттерн по \citet{karpman1941}, эмпирически воспроизведённый у нас на $n=144$); может иметь оба одновременно (вторичный паттерн, $n=65$, $\mathrm{ID} \approx +0.90$ и $\mathrm{CA} \approx +0.98$). Их независимость --- один из главных эмпирических результатов статьи (§\ref{subsec:architecture}), и из неё вытекают существенно различающиеся клинические стратегии.

Эта инверсия рамки поддерживается двумя внешними эмпирическими опорами, которые мы считаем наиболее сильными аргументами в её пользу. Во-первых, метааналитическая оценка \citet{depow2025}: трейт-показатели эмпатии объясняют всего 3--15\% дисперсии эмпатических переживаний в повседневной жизни. Если бы эмпатия была стабильной индивидуальной чертой, которая некоторым свойственна, а некоторым нет, стандартные трейт-инструменты улавливали бы значительно бóльшую долю дисперсии. Фактические 3--15\% указывают на то, что эмпатический отклик --- это state, сильно зависящий от ситуативных факторов, среди которых центральную роль играет \emph{активация или деактивация блокирующих механизмов}, а не колебание базового ресурса. Во-вторых, результат \citet{ogawa1997}: отсутствие родительского отклика в младенчестве предсказывает диссоциативную симптоматику 19 лет спустя с эффектом размера около 50\% объяснённой дисперсии --- сильнее, чем предсказывают сами травматические события. Это согласуется с гипотезой \citet{fonagy2002}: не стрессовое событие как таковое, а отсутствие ментализующего поля, в котором событие могло бы быть переработано, формирует структурную диссоциативную защиту. То есть механизм блокировки формируется именно там, где разрушается или никогда не складывается способность к со-переживанию взрослым внутренних состояний ребёнка.

Третья опора --- структурная. \citet{digiuseppe2024} в сетевом анализе защитных механизмов показали, что 16 различных защит формируют плотно связанную интегрированную систему (не независимые черты), в которой самоутверждение является центральным узлом, а остальные защиты --- периферические узлы, тянущие на себя разные функциональные нагрузки. В наших данных корреляции между подшкалами пассивной блокировки (диссоциация, подавление, аффективная изоляция, притупление, внутренний диссонанс, маскирование) лежат в диапазоне $r = 0.77$--$0.94$. Это не артефакт избыточности инструмента, а воспроизведение того, что сетевой анализ Di Giuseppe обнаружил независимо: защитные механизмы суть \emph{одна интегрированная функциональная система}, а не набор параллельных индивидуальных черт. Именно поэтому в нашей 6-factor CFA эти подшкалы сливаются в единый латентный фактор ID --- они психометрически не дискриминантны не потому, что инструмент их не различает, а потому, что в реальности они --- разные феноменологические выражения одной системы подавления.

Эти три опоры --- низкая объяснительная сила трейт-эмпатии, мощное предсказание диссоциации через отсутствие ментализующего отклика, и сетевая интегрированность защит --- независимо сходятся на одном выводе: центральный объект измерения --- \emph{механизмы не-эмпатии}, а не эмпатия как черта. Модель \DEIC{} построена на этом переворачивании, и все её последующие утверждения следует читать через эту рамку.

\subsection{Пять точек, в которых традиции касались друг друга без встречи}

Причина, по которой \DEIC{}-архитектура имеет шанс быть полезной, не в амбициозности, а в узости. Мы выбрали для моделирования те конструкты, в которых разные традиции эмпирически описывают \emph{очень похожие} вещи, но не имели общего инструмента, чтобы это подтвердить.

\paragraph{Внутренний диссонанс (Internal Dissonance, ID).}
В психоаналитической традиции --- это расщепление и отщепление эго \citep{ferenczi1933, klein1946}; в теории привязанности --- дезорганизованные внутренние рабочие модели \citep{main1986}; в клинической диссоциативной литературе --- фрагментация в ответ на травму \citep{putnam1997, vanderkolk2014}. Все три описывают то же самое явление: внутренние состояния человека несут признак активной блокировки, регуляции, защитной работы, а не просто содержания. Психометрически это должно быть одним фактором, если инструмент правильно сконструирован. Мы покажем, что это так.

\paragraph{Детское отчуждение (Childhood Alienation, CA).}
В теории привязанности --- отсутствие secure base \citep{bowlby1969, ainsworth1978}; в mentalization theory --- отсутствие ментализующего отклика \citep{fonagy2002}; в клинической ACE-литературе --- стрессовые события детства \citep{felitti1998}. CA --- это средовое условие, в котором внутренние состояния ребёнка не встречают адекватного отклика. Мы его узко инструментируем (не полный ACE-индекс), потому что, согласно гипотезам Fonagy, именно \emph{отсутствие ментализующего поля}, а не сам факт стрессовых событий, предиктирует последующую дисрегуляцию. Содержательный смысл этой узкой формулировки --- ребёнок учится различать и удерживать собственные состояния через взрослого, который способен эти состояния увидеть и назвать. Если такого взрослого рядом нет, ребёнок не теряет «информацию о других» отдельно от «информации о себе»: он не развивает общую способность метаболизировать аффективный опыт как таковой. Поэтому CA $\to$ ID в нашей модели --- это не «травма снизила эмпатию к окружающим»; это «отсутствие ментализующего зеркала не позволило сформироваться способности выдерживать аффективный материал в обе стороны сразу». Клиническая сигнатура взрослого с высоким CA и высоким ID --- одинаковая непроницаемость и наружу, и вовнутрь: такой человек и к собственной грусти не проходит, и к грусти партнёра не проходит, потому что единственный доступный механизм --- автоматическая общая блокировка, а не селективное закрытие одного из направлений.

\paragraph{Психопатия как моральная подпись (P).}
Клиническая традиция \citep{cleckley1941, hare1991} описывает психопатию в рамке аффективного дефицита и инструментального использования других; нейронаука Блэра \citep{blair1995, blair2013} --- через сниженную амигдальную реакцию на чужой дистресс; эволюционная психология \citep{mealey1995, lalumiere2001} --- как стабильную на низкой частоте стратегию. Мы намеренно не воспроизводим дефицитарную рамку. В \DEIC{} психопатический паттерн расщепляется на два слоя: \textbf{архитектурный} (условное, не безусловное, аффективное сцепление --- морально нейтральная вариация, функциональная в ряде социальных ролей, см. §\ref{sec:discussion}.5) и \textbf{нарративный} (P --- устойчивое утверждение «это разрешено для меня», определяющее, как архитектура используется социально). Фактор P моделирует именно второй слой. Этот двухслойный сдвиг --- от «сломанной функции» к «архитектуре плюс нарратив» --- меняет клинические, юридические и политические следствия и является одним из ключевых вкладов работы.

\paragraph{Когнитивная vs аффективная эмпатия (\Ecog{} vs \EAF{}).}
В нейропсихологии --- две системы с различимыми субстратами \citep{shamaytsoory2011, decety2011}; в теории mentalization --- ментализация vs embodied empathy \citep{fonagy2002}; в клинической работе --- ключевое различение для понимания того, почему «умный психопат» может быть точен в предсказаниях поведения, но безразличен к страданию. Мы показываем эмпирически, что эти два конструкта разделимы и по-разному связаны с P. Повседневная иллюстрация этого различения --- одно из самых структурно-информативных наблюдений, которое обычно игнорируется в trait-измерительной литературе: один и тот же человек может плакать над судьбой литературного или киногероя и оставаться холодным к страданию конкретного реального человека рядом с ним. Это не противоречие и не «лицемерие» --- это прямая сигнатура расщепления каналов. Идентификация с нарративом (вымышленный герой внутри безопасной рамки повествования) задействует преимущественно \Ecog{}-линию: ментализация точки зрения без автоматического аффективного сцепления с реальным, угрожающим, небезопасным присутствием другого. Реальный страдающий человек требует \EAF{} в её полной форме --- аффективного резонанса в ситуации, которая не управляется наррартивной рамкой и может предъявить требования. Когда \EAF{} заблокирована защитной структурой (ID) или регулируется условно (P), канал через \Ecog{} остаётся частично или полностью сохранным, и человек переживает состояние, в котором вымышленное страдание трогает, а реальное --- нет. В DEIC-архитектуре это не аномалия, а предсказуемое следствие разделимости двух каналов.

\paragraph{Внимание как witnessing (A).}
В контемплативной традиции --- \textit{satipaṭṭhāna} (буддизм), \textit{sākṣī} (адвайта-веданта), \textit{rigpa} (дзогчен), «самовоспоминание» \citep[Гурджиев через][]{ouspensky1949}; в психоаналитической традиции --- наблюдающее Я \citep{sterba1934, kris1956}; в современной mindfulness-науке --- meta-awareness \citep{dunne2011, josipovic2014}; в mentalization theory --- рефлективная функция \citep{fonagy1997}. Во всех этих линиях описывается структурно особый слой сознания: не содержание психики, а \emph{отношение наблюдения} к собственному содержанию. Мы будем возвращаться к A многократно --- это самый теоретически богатый и одновременно самый сложный для инструментирования фактор в модели.

\medskip
Важно: мы не претендуем на то, что \DEIC{} описывает эти традиции \emph{целиком} или \emph{первыми} ловит структуру, которую они описывают. Каждая из перечисленных литератур измеряла свой конструкт с большей глубиной, чем мы. Наш вклад --- в том, что мы измерили \emph{все перечисленные конструкты одной выборкой, одним инструментом, с общими методологическими стандартами}. Это позволяет впервые протестировать их взаимные отношения не на уровне теоретической реконструкции, а на уровне коэффициентов путевой модели.

\subsection{Балканизация как failure mode, а не дисциплина}

Прежде чем переходить к методу, сформулируем позицию открыто. Традиция считать, что «узкая рамка = строгая наука», в психологии XX века обернулась тем, что несовместимые микротеории сосуществуют как норма и называются дисциплиной. Теория нарциссизма не согласована с теорией психопатии; обе не согласованы с теорией привязанности; когнитивная эмпатия измеряется инструментами, архитектурно несовместимыми с аффективной. Это именно то, что \citet{meehl1978} диагностировал полвека назад как failure mode психологического теоретизирования. Физика, химия и биология требуют когерентности через масштабы; психология позволила себе отказаться от этого требования под риторикой «осторожности».

Мы не занимаем эту осторожность как добродетель. Мы утверждаем: если в данных есть конструкты, которые разные традиции описывают как одно и то же явление под разными именами, психометрия должна это продемонстрировать. Если \DEIC{}-архитектура делает предсказания, она обязана их сформулировать фальсифицируемо --- именно так, а не через «узкое ограничение scope'а», которое выводит из-под проверки то, что на самом деле модель предсказывает. Это методологическая ставка статьи: амбиция по содержанию, дисциплина по тому, как ставится вопрос.

\subsection{Позиционирование: три слоя, одно предложение}

Относительно трёх референтных плоскостей работа располагается так. К \textbf{контемплативной традиции} (буддизм, адвайта, дзогчен, исихазм, Гурджиев, психоанализ у своего созерцательного края) \DEIC{} стоит в позиции \emph{ученика-переводчика}: конструкты, которые эти традиции различали без психометрических инструментов тысячелетиями, у нас получают координаты --- и наш вклад именно в этом, не в «новом интегрировании». К \textbf{академической психологии} \DEIC{} стоит в \emph{конфронтационной} позиции: балканизация --- не дисциплина, а патология поля, и каркас общих референтов для того, что школы описывают в несовместимых языках, --- это то, что поле должно иметь, но не имеет. К \textbf{интегральной традиции} (Уилбер и AQAL, structural precedents) \DEIC{} стоит в \emph{структурно-обратной} позиции: не master-map, на которую традиции попадают как частные случаи, а операциональный диалект тех же традиций, развёрнутый против фрагментации.

Одно предложение: \DEIC{} --- операциональный диалект того, что контемплативная традиция измеряла без инструментов, развёрнутый против балканизации академической психологии. Ученики традиции, вызов полю.

\subsection{Что \DEIC{} не обещает}

Поскольку программа позиционируется как каркас, пересекающий несколько традиций, мы сразу очерчиваем границы. \DEIC{} \textbf{не} претендует:
\begin{itemize}
    \item Быть первым или единственным унифицирующим фреймворком эмпатии. Унификация --- не наш глагол.
    \item Содержать окончательную операционализацию A, witnessing, или mindfulness. A в текущем виде --- формативный композит, не рефлексивный латент; это ограничение эмпирически установлено (§\ref{sec:results}.7) и исправляется в следующей волне.
    \item Заменять любую из традиций, которые мы цитируем. Психоанализ, теория привязанности, когнитивная нейронаука эмпатии, зеркально-нейронная литература, клиническая традиция психопатии, созерцательные школы --- каждая в своей парадигме имеет глубину и специфичность, которую мы не воспроизводим.
    \item Предлагать теорию сознания или объяснение hard problem \citep{chalmers1995}. Мы измеряем \emph{функции} внутри сознания; мы ничего не утверждаем о природе субъективного опыта как такового.
    \item Давать clinical recommendations на основе одной выборки. Всё клинически-prescriptive в Discussion сформулировано как \emph{гипотезы для тестирования}, не как немедленные протоколы.
\end{itemize}

Эти ограничения --- не reluctant retreat, а структурное условие возможности сделать работу, которая не превратится в мета-теоретическую претензию. Программа строится из многих узких, фальсифицируемых работ; данная статья --- одна из них.

\subsection{Структура статьи}

В разделе~\ref{sec:methods} мы описываем инструмент DEIC-Inventory, операционализацию шести факторов через 6-factor oblique CFA, конструкцию A как формативного композита, и статистический подход (conditional process SEM, параллельную медиацию, unified SEM, A-CFA). В разделе~\ref{sec:results} мы представляем факторную структуру, partial correlations, центральный conditional-process анализ, параллельную медиацию, unified SEM со всеми путями одновременно, эмпирическое подтверждение интегрированного выжившего, топологическое различение первичной и вторичной психопатии, и эмпирическую демонстрацию формативной природы A. В разделе~\ref{sec:discussion} мы обсуждаем импликации архитектуры, выносим A как онтологически асимметричный фактор (field-condition, not field-content), формулируем фальсифицируемые предсказания и точки расхождения с существующими теориями, и указываем место статьи в будущей программе. Ограничения и направления будущей работы --- в разделах~\ref{sec:limitations} и~\ref{sec:future}.
